\chapter{Термодинамическая ресурсная граница автономных квантовых часов}
\label{ch:v8-bounded-strength-autonomous-clock-thermodynamic-limit-gate}

Предыдущий гейт дал точное конечное исполнение, но часы не возвращались в
начало, а локальная цепь истории требовала ресурса, растущего с объёмом.
Это ещё не общий запрет: стационарные локальные гамильтонианы способны
автономно исполнять запрограммированные последовательности, а конечные
квазиидеальные часы подавляют обратное воздействие экспоненциально по своей
размерности. Поэтому проверим не существование вообще, а ресурс, необходимый
для равномерной точности.

\section{Конечный объём}

Пусть однократное управляемое действие часов размерности $d$ приближает
идеальное локальное действие с оценкой
\begin{equation}
 \varepsilon_d\leq A e^{-cd},\qquad A,c>0.
 \label{eq:v8-quasi-ideal-clock-error}
\end{equation}
Для двухслойного конвейера на кольце длины $L$ требуется $N_L=2L$
локальных действий. Телескопическая оценка для композиции каналов даёт
\begin{equation}
 \left\|\widetilde\Phi_L-\Phi_L\right\|_\diamond
 \leq N_L\varepsilon_d
 \leq 2LAe^{-cd}.
 \label{eq:v8-clock-global-error}
\end{equation}
Следовательно, для любого конечного $L$ и заданного $0<\delta<1$ достаточно
выбрать
\begin{equation}
 d\geq \frac1c\log\frac{2AL}{\delta}.
 \label{eq:v8-clock-dimension-schedule}
\end{equation}
Если средняя энергия квазиидеальных часов растёт не быстрее $E_0d$, то
достаточный энергетический ресурс растёт как
\begin{equation}
 E_d=O\!\left(\log\frac{L}{\delta}\right).
 \label{eq:v8-clock-energy-schedule}
\end{equation}

\begin{theorem}[Конечномерное автономное приближение]
При оценке \eqref{eq:v8-quasi-ideal-clock-error} любой конечный
двухслойный конвейер допускает автономное приближение с произвольно малой
заданной ошибкой. Достаточная размерность часов растёт лишь логарифмически
с объёмом и обратной точностью.
\end{theorem}

\section{Два неэквивалентных термодинамических предела}

При фиксированном $d$ оценка \eqref{eq:v8-clock-global-error} не равномерна
по $L$. Поэтому данная конструкция не доказывает сходимость полного канала
на всей бесконечной цепи в норме $\|\cdot\|_\diamond$. Это граница
сертификата, а не общий запрет всех взаимодействующих часовых моделей.

Для наблюдаемой с конечным носителем при фиксированном времени ситуация
иная. Если локальный гамильтониан имеет конечную скорость распространения,
на наблюдаемую влияет лишь конечное число $N_{\rm loc}$ действий внутри её
причинного конуса. Тогда
\begin{equation}
 \varepsilon_{\rm loc}\leq N_{\rm loc}Ae^{-cd},
 \label{eq:v8-clock-local-observable-error}
\end{equation}
где $N_{\rm loc}$ не зависит от полного объёма $L$. Тем самым локальный
термодинамический предел допускается при фиксированных конечных часах,
хотя глобальная норма канала остаётся незамкнутой.

\begin{nogocriterion}[Запрет отождествления локального и глобального пределов]
Нельзя из конечномерной точности часов заключать равномерную точность
полного бесконечного конвейера. Нельзя также объявлять этот результат
общим запретом автономных часов: он утверждает только, что имеющийся
экспоненциальный сертификат требует роста $d$ для глобального контроля.
Точный возврат часов и абсолютная длительность такта по-прежнему не
выведены.
\end{nogocriterion}

\begin{protocol}
\begin{enumerate}
 \item ошибка одного управления: \textbf{$\varepsilon_d\leq Ae^{-cd}$};
 \item число действий на кольце: \textbf{$2L$};
 \item глобальная оценка: \textbf{$2LAe^{-cd}$};
 \item достаточная размерность: \textbf{$d\geq c^{-1}\log(2AL/\delta)$};
 \item достаточный энергетический рост: \textbf{$O(\log(L/\delta))$};
 \item точный конечномерный конвейер: \textbf{не получен};
 \item глобальная равномерность при фиксированном $d$: \textbf{не доказана};
 \item локальный предел конечных наблюдаемых: \textbf{допущен условно};
 \item абсолютная длительность такта: \textbf{не выведена};
 \item следующий гейт: \textbf{локальный предел часового квантового
 марковского процесса и размерная привязка времени}.
\end{enumerate}
\end{protocol}

Машинный сертификат сохранён в
\path{s2t_v8_bounded_strength_autonomous_clock_thermodynamic_limit_gate_results.json}.

\begin{thebibliography}{9}
\bibitem{v8-bounded-clock-woods}
M.~P.~Woods, R.~Silva, J.~Oppenheim,
\emph{Autonomous quantum machines and the finite sized Quasi-Ideal clock},
Ann. Henri Poincar\'e \textbf{20} (2019), 125--218;
arXiv:1607.04591.
\bibitem{v8-bounded-clock-chase}
B.~A.~Chase, A.~J.~Landahl,
\emph{Universal quantum walks and adiabatic algorithms by 1D Hamiltonians},
arXiv:0802.1207.
\end{thebibliography}